% !TEX root = ../matan.tex

\section{Двойственные пространства, комплексные пространства и ряды функций}

\subsection{Двойственное пространство}

\begin{Definition}{Линейный функционал (напоминание)}{}
	Пусть $X$ --- нормированное линейное пространство над полем $\mathbb K$, где $\mathbb K=\R$ или $\mathbb K=\C$.
	
	Отображение $\mathcal A\colon X\to \mathbb K$ называется линейным функционалом, если
	\[
	\mathcal A(\lambda x+y)=\lambda \mathcal A(x)+\mathcal A(y)
	\]
	для всех $x,y\in X$ и $\lambda\in\mathbb K$.
	
	Если $\mathcal A$ непрерывен, то он называется непрерывным линейным функционалом.
\end{Definition}

\begin{Definition}{Двойственное пространство}{}
	Двойственным пространством к $X$ называется пространство всех непрерывных линейных функционалов на $X$:
	\[
	X^*=\{\mathcal A\colon X\to\mathbb K \mid \mathcal A \text{ линейный и непрерывный}\}.
	\]
	
	На $X^*$ вводится норма
	\[
	\norma{\mathcal A}_{X^*}=\sup_{\norma{x}_X\le 1}\abs{\mathcal A(x)}.
	\]
\end{Definition}

\begin{Remark}{}{}
	Элементы $X$ --- это векторы, а элементы $X^*$ --- это непрерывные линейные способы получить из вектора число. Поэтому $X^*$ часто называют пространством линейных измерений на $X$.
\end{Remark}

\begin{Theorem}{Теорема Рисса}{}
	Пусть $\mathcal H$ --- гильбертово пространство, а $\mathcal A\in \mathcal H^*$ --- непрерывный линейный функционал. Тогда существует единственный вектор $y\in \mathcal H$ такой, что
	\[
	\mathcal A(x)=\scal{x}{y}_{\mathcal H}
	\quad \text{для всех } x\in \mathcal H.
	\]
	
	Иначе говоря, всякий непрерывный линейный функционал на гильбертовом пространстве задаётся скалярным произведением с некоторым фиксированным вектором.
\end{Theorem}

\begin{Remark}{Отождествление $\mathcal H$ и $\mathcal H^*$}{}
	Теорема Рисса позволяет отождествлять гильбертово пространство со своим двойственным пространством:
	\[
	\mathcal H^*\simeq \mathcal H.
	\]
	Например, $(\ell^2)^*\simeq \ell^2$.
\end{Remark}

\begin{Example}{Пространства $\ell^p$ и их двойственные пространства}{}
	Для $1\le p<\infty$ пространство $\ell^p$ состоит из всех последовательностей $x=(x_1,x_2,\ldots)$, где $x_k\in\mathbb K$, таких что
	\[
	\norma{x}_{\ell^p}
	=
	\left(\sum_{k=1}^{\infty}\abs{x_k}^p\right)^{1/p}
	<\infty.
	\]
	
	Пространство $\ell^p$ является банаховым при всех $1\le p<\infty$. Гильбертовым оно является только при $p=2$.
	
	Если $1<p<\infty$ и число $q$ задаётся условием $\frac1p+\frac1q=1$, то
	\[
	(\ell^p)^*\simeq \ell^q.
	\]
	Точнее, для всякого $\mathcal A\in(\ell^p)^*$ существует единственный $y=(y_k)\in\ell^q$ такой, что
	\[
	\mathcal A(x)=\sum_{k=1}^{\infty}x_k y_k
	\quad \text{для всех } x=(x_k)\in\ell^p.
	\]
	
	Корректность этой формулы обеспечивается неравенством Гёльдера:
	\[
	\abs{\sum_{k=1}^{\infty}x_k y_k}
	\le
	\left(\sum_{k=1}^{\infty}\abs{x_k}^p\right)^{1/p}
	\left(\sum_{k=1}^{\infty}\abs{y_k}^q\right)^{1/q}
	=
	\norma{x}_{\ell^p}\norma{y}_{\ell^q}.
	\]
	Следовательно, ряд $\sum x_k y_k$ сходится абсолютно, а соответствующий функционал непрерывен.
\end{Example}

\begin{Definition}{Пространство $C(K)$}{}
	Пусть $K$ --- компактное хаусдорфово пространство; в частности, можно думать о компакте $K\subset\R^n$. Обозначим через $C(K)$ пространство непрерывных функций $f\colon K\to\mathbb K$.
	
	На $C(K)$ вводится равномерная норма:
	\[
	\norma{f}_{C(K)}
	=
	\norma{f}_{\infty}
	=
	\sup_{x\in K}\abs{f(x)}.
	\]
	С этой нормой $C(K)$ является банаховым пространством.
\end{Definition}

\subsection{Комплексные числа}

\begin{Definition}{Поле комплексных чисел}{}
	Комплексные числа можно рассматривать как множество $\R^2$ с операциями сложения и умножения. Элемент $z\in\C$ записывается в виде $z=x+iy$, где $x,y\in\R$, $i=(0,1)$, $1=(1,0)$ и $i^2=-1$.
\end{Definition}

Если $z_1=x_1+iy_1$ и $z_2=x_2+iy_2$, то
\[
z_1z_2
=
(x_1x_2-y_1y_2)
+
i(x_1y_2+x_2y_1).
\]

\begin{Definition}{Сопряжение, модуль и обратный элемент}{}
	Для $z=x+iy$ определяются:
	\[
	\overline z=x-iy,
	\qquad
	\abs z=\sqrt{x^2+y^2},
	\qquad
	z\overline z=\abs z^2.
	\]
	Если $z\ne0$, то
	\[
	\frac1z
	=
	\frac{\overline z}{\abs z^2}
	=
	\frac{x-iy}{x^2+y^2}
	=
	\frac{x}{x^2+y^2}
	-
	i\frac{y}{x^2+y^2}.
	\]
\end{Definition}

\begin{Definition}{Аргумент комплексного числа}{}
	Если $z=x+iy\ne0$, то аргументом $z$ называется число $\varphi$ такое, что
	\[
	x=\abs z\cos\varphi,
	\qquad
	y=\abs z\sin\varphi.
	\]
	Пишут $\varphi=\arg z$. Аргумент определён с точностью до прибавления $2\pi k$, $k\in\mathbb Z$.
\end{Definition}

\begin{Statement}{Формула Эйлера}{}
	Для любого $\varphi\in\R$
	\[
	e^{i\varphi}=\cos\varphi+i\sin\varphi.
	\]
	Отсюда
	\[
	\cos\varphi=\frac{e^{i\varphi}+e^{-i\varphi}}2,
	\qquad
	\sin\varphi=\frac{e^{i\varphi}-e^{-i\varphi}}{2i}.
	\]
\end{Statement}

\begin{Remark}{Полярная форма}{}
	Ненулевое комплексное число можно записать в виде
	\[
	z=\abs z e^{i\arg z}.
	\]
	Если $z_1,z_2\ne0$, то
	\[
	z_1z_2
	=
	\abs{z_1}\abs{z_2}
	e^{i(\arg z_1+\arg z_2)}.
	\]
	То есть при умножении комплексных чисел модули перемножаются, а аргументы складываются.
\end{Remark}

\subsection{Полилинейные отображения}

\begin{Definition}{Полилинейное отображение}{}
	Пусть $X_1,\ldots,X_N,Y$ --- линейные пространства. Отображение
	\[
	\mathcal A\colon X_1\times\cdots\times X_N\to Y
	\]
	называется полилинейным, если оно линейно по каждому аргументу при фиксированных остальных аргументах.
	
	То есть для каждого $k=1,\ldots,N$:
	\[
	\mathcal A(x_1,\ldots,x_{k-1},\lambda x_k+y_k,x_{k+1},\ldots,x_N)
	=
	\lambda\mathcal A(x_1,\ldots,x_k,\ldots,x_N)
	+
	\mathcal A(x_1,\ldots,y_k,\ldots,x_N).
	\]
\end{Definition}

\begin{Remark}{Прямое произведение пространств}{}
	Элемент пространства $X_1\times\cdots\times X_N$ имеет вид $x=(x_1,\ldots,x_N)$, где $x_k\in X_k$. Если все пространства совпадают с одним пространством $X$, то пишут $X^N$.
\end{Remark}

\subsection{Комплексные пространства со скалярным произведением}

\begin{Definition}{Скалярное произведение над $\C$}{}
	Пусть $\mathcal H$ --- линейное пространство над $\C$. Скалярным произведением на $\mathcal H$ называется отображение $\scal{\cdot}{\cdot}\colon \mathcal H\times\mathcal H\to\C$ со свойствами:
	
	\begin{enumerate}
		\item $\scal{x}{x}\in\R$, $\scal{x}{x}\ge0$, причём $\scal{x}{x}=0\iff x=0$;
		\item $\scal{\lambda x}{y}=\lambda\scal{x}{y}$;
		\item $\scal{x}{y}=\overline{\scal{y}{x}}$, поэтому $\scal{x}{\lambda y}=\overline\lambda\scal{x}{y}$;
		\item $\scal{x+y}{z}=\scal{x}{z}+\scal{y}{z}$.
	\end{enumerate}
	
	Здесь используется соглашение: скалярное произведение линейно по первому аргументу и сопряжённо-линейно по второму.
\end{Definition}

\begin{Example}{$\ell^2$ над $\C$}{}
	В комплексном пространстве $\ell^2$ скалярное произведение задаётся формулой
	\[
	\scal{x}{y}_{\ell^2}
	=
	\sum_{k=1}^{\infty}x_k\overline{y_k}.
	\]
	Соответствующая норма:
	\[
	\norma{x}_{\ell^2}
	=
	\left(\sum_{k=1}^{\infty}\abs{x_k}^2\right)^{1/2}.
	\]
\end{Example}

\subsection{Последовательности и ряды функций}

\begin{Definition}{Последовательность функций}{}
	Пусть $\Omega$ --- множество, а $f_k\colon\Omega\to\C$. Тогда $\{f_k\}_{k\in\N}$ называется последовательностью функций.
	
	Для каждой фиксированной точки $z_0\in\Omega$ получается числовая последовательность $\{f_k(z_0)\}_{k\in\N}$.
\end{Definition}

\begin{Definition}{Поточечная сходимость последовательности функций}{}
	Последовательность функций $f_k\colon\Omega\to\C$ сходится поточечно к функции $f\colon\Omega\to\C$, если для каждого $x\in\Omega$ числовая последовательность $f_k(x)$ сходится к $f(x)$.
	
	То есть:
	\[
	\forall x\in\Omega:
	\quad
	\lim_{k\to\infty}f_k(x)=f(x).
	\]
\end{Definition}

\begin{Definition}{Функциональный ряд}{}
	Пусть $f_k\colon\Omega\to\C$. Формальное выражение
	\[
	\sum_{k=1}^{\infty}f_k
	\]
	называется функциональным рядом.
	
	Его частичные суммы:
	\[
	S_n(x)=\sum_{k=1}^{n}f_k(x).
	\]
	Ряд сходится поточечно к функции $S$, если $S_n(x)\to S(x)$ для каждого $x\in\Omega$.
\end{Definition}

\begin{Example}{Показательная функция и тригонометрические функции}{}
	Для любого $z\in\C$ ряд $\sum_{k=0}^{\infty}\frac{z^k}{k!}$ сходится абсолютно. Его сумма обозначается $e^z$:
	\[
	e^z=\sum_{k=0}^{\infty}\frac{z^k}{k!}.
	\]
	
	Также:
	\[
	\cos z=\sum_{k=0}^{\infty}(-1)^k\frac{z^{2k}}{(2k)!},
	\qquad
	\sin z=\sum_{k=0}^{\infty}(-1)^k\frac{z^{2k+1}}{(2k+1)!}.
	\]
\end{Example}

\subsection{Равномерная сходимость}

\begin{Definition}{Равномерная сходимость последовательности функций}{}
	Пусть $f_n\colon\Omega\to\C$. Последовательность $\{f_n\}$ равномерно сходится к функции $f\colon\Omega\to\C$, если
	\[
	\forall\eps>0\ \exists N\ \forall n>N\ \forall x\in\Omega:
	\abs{f_n(x)-f(x)}<\eps.
	\]
	Эквивалентно:
	\[
	\norma{f_n-f}_{\infty}
	=
	\sup_{x\in\Omega}\abs{f_n(x)-f(x)}
	\to0.
	\]
\end{Definition}

\begin{Statement}{Критерий Коши для равномерной сходимости}{}
	Последовательность $\{f_n\}$ равномерно сходится на $\Omega$ тогда и только тогда, когда
	\[
	\forall\eps>0\ \exists N\ \forall n,m>N:
	\sup_{x\in\Omega}\abs{f_n(x)-f_m(x)}<\eps.
	\]
	Иными словами, $\{f_n\}$ является фундаментальной последовательностью в норме $\norma{\cdot}_{\infty}$.
\end{Statement}

\begin{Statement}{Равномерная сходимость влечёт поточечную}{}
	Если $f_n\to f$ равномерно на $\Omega$, то $f_n(x)\to f(x)$ для каждого $x\in\Omega$.
\end{Statement}

\begin{proof}
	Из равномерной сходимости для любого $\eps>0$ существует $N$ такое, что при $n>N$ выполнено $\sup_{x\in\Omega}\abs{f_n(x)-f(x)}<\eps$. Тогда для каждого фиксированного $x\in\Omega$ имеем $\abs{f_n(x)-f(x)}<\eps$. Значит, $f_n(x)\to f(x)$.
\end{proof}

\begin{Remark}{Обратное неверно}{}
	Поточечная сходимость вообще говоря не влечёт равномерную. Например, $f_n(x)=x^n$ на $(0,1)$ поточечно сходится к нулю, но не равномерно.
\end{Remark}

\begin{Definition}{Равномерная ограниченность}{}
	Последовательность функций $\{f_k\}$ называется равномерно ограниченной на $\Omega$, если существует $M>0$ такое, что
	\[
	\sup_{x\in\Omega}\abs{f_k(x)}\le M
	\quad \text{для всех } k.
	\]
\end{Definition}

\begin{Theorem}{}{}
	Пусть $I\subset\R$ --- промежуток, $f_n\colon I\to\C$, все $f_n$ непрерывны на $I$, и $f_n\to f$ равномерно на $I$. Тогда $f$ непрерывна на $I$.
\end{Theorem}

\begin{proof}
	Докажем непрерывность $f$ в произвольной точке $x_0\in I$. Пусть $\eps>0$. По равномерной сходимости существует $N$ такое, что для всех $x\in I$:
	\[
	\abs{f_N(x)-f(x)}<\frac{\eps}{3}.
	\]
	Так как $f_N$ непрерывна в точке $x_0$, существует $\delta>0$ такое, что при $x\in I$ и $\abs{x-x_0}<\delta$ выполнено
	\[
	\abs{f_N(x)-f_N(x_0)}<\frac{\eps}{3}.
	\]
	Тогда при $\abs{x-x_0}<\delta$:
	\[
	\abs{f(x)-f(x_0)}
	\le
	\abs{f(x)-f_N(x)}
	+
	\abs{f_N(x)-f_N(x_0)}
	+
	\abs{f_N(x_0)-f(x_0)}
	<\eps.
	\]
	Значит, $f$ непрерывна в $x_0$.
\end{proof}

\begin{Statement}{}{}
	Пусть $\{f_k\}$ равномерно ограничена на $\Omega$, а $g_k\to0$ равномерно на $\Omega$. Тогда $f_k g_k\to0$ равномерно на $\Omega$.
\end{Statement}

\begin{proof}
	Пусть $\abs{f_k(x)}\le M$ для всех $k$ и $x\in\Omega$. Тогда
	\[
	\sup_{x\in\Omega}\abs{f_k(x)g_k(x)}
	\le
	M\sup_{x\in\Omega}\abs{g_k(x)}.
	\]
	Так как $g_k\to0$ равномерно, правая часть стремится к нулю.
\end{proof}

\subsection{Равномерная сходимость рядов функций}

\begin{Definition}{Равномерная сходимость функционального ряда}{}
	Ряд $\sum_{k=1}^{\infty}f_k$ равномерно сходится на $\Omega$, если его частичные суммы $S_n(x)=\sum_{k=1}^{n}f_k(x)$ равномерно сходятся на $\Omega$.
\end{Definition}

\begin{Statement}{Критерий Коши для рядов функций}{}
	Ряд $\sum_{k=1}^{\infty}f_k$ равномерно сходится на $\Omega$ тогда и только тогда, когда
	\[
	\forall\eps>0\ \exists N\ \forall m>n>N:
	\sup_{x\in\Omega}
	\abs{\sum_{k=n+1}^{m}f_k(x)}
	<\eps.
	\]
\end{Statement}

\begin{Statement}{Необходимое условие равномерной сходимости ряда}{}
	Если ряд $\sum_{k=1}^{\infty}f_k$ равномерно сходится на $\Omega$, то $f_k\to0$ равномерно на $\Omega$.
\end{Statement}

\begin{proof}
	По критерию Коши для рядов при любом $\eps>0$ существует $N$ такое, что при $m>n>N$:
	\[
	\sup_{x\in\Omega}
	\abs{\sum_{k=n+1}^{m}f_k(x)}
	<\eps.
	\]
	Берём $m=n+1$. Тогда $\sup_{x\in\Omega}\abs{f_{n+1}(x)}<\eps$. Значит, $\norma{f_k}_{\infty}\to0$.
\end{proof}

\begin{Statement}{Удобная формулировка равномерной малости}{}
	Последовательность $f_k\colon\Omega\to\C$ равномерно сходится к нулю тогда и только тогда, когда существует числовая последовательность $\eps_k\to0$ такая, что
	\[
	\sup_{x\in\Omega}\abs{f_k(x)}\le \eps_k
	\quad \text{для всех } k.
	\]
\end{Statement}

\begin{Theorem}{Признак Вейерштрасса}{}
	Пусть $f_k\colon\Omega\to\C$. Если существуют числа $a_k\ge0$ такие, что $\abs{f_k(x)}\le a_k$ для всех $x\in\Omega$, и числовой ряд $\sum_{k=1}^{\infty}a_k$ сходится, то функциональный ряд $\sum_{k=1}^{\infty}f_k$ сходится равномерно на $\Omega$.
\end{Theorem}

\begin{proof}
	Так как $\sum a_k$ сходится, то по критерию Коши для числовых рядов для любого $\eps>0$ существует $N$ такое, что при $m>n>N$:
	\[
	\sum_{k=n+1}^{m}a_k<\eps.
	\]
	Тогда для всех $x\in\Omega$:
	\[
	\abs{\sum_{k=n+1}^{m}f_k(x)}
	\le
	\sum_{k=n+1}^{m}\abs{f_k(x)}
	\le
	\sum_{k=n+1}^{m}a_k
	<\eps.
	\]
	Следовательно,
	\[
	\sup_{x\in\Omega}
	\abs{\sum_{k=n+1}^{m}f_k(x)}
	<\eps.
	\]
	По критерию Коши ряд $\sum f_k$ сходится равномерно.
\end{proof}

\subsection{Признаки Абеля и Дирихле}

\begin{Theorem}{Признак Абеля}{}
	Пусть $f_k,g_k\colon\Omega\to\C$. Предположим, что:
	
	\begin{enumerate}
		\item ряд $\sum_{k=1}^{\infty}f_k$ сходится равномерно на $\Omega$;
		\item последовательность $\{g_k\}$ равномерно ограничена: $\abs{g_k(x)}\le M$ для всех $k$ и всех $x\in\Omega$;
		\item для каждого $x\in\Omega$ последовательность $\{g_k(x)\}$ монотонна.
	\end{enumerate}
	
	Тогда ряд $\sum_{k=1}^{\infty}f_k(x)g_k(x)$ сходится равномерно на $\Omega$.
\end{Theorem}

\begin{proof}
	Докажем по критерию Коши. Пусть $\eps>0$. Так как $\sum f_k$ сходится равномерно, то существует $N$ такое, что при любых $m>n>N$:
	\[
	\sup_{x\in\Omega}
	\abs{\sum_{k=n+1}^{m}f_k(x)}
	<\eps.
	\]
	
	Используем преобразование Абеля. Для фиксированного $x\in\Omega$ обозначим
	\[
	F_j(x)=\sum_{k=n+1}^{j}f_k(x).
	\]
	Тогда
	\[
	\sum_{k=n+1}^{m}f_k(x)g_k(x)
	=
	F_m(x)g_m(x)
	+
	\sum_{k=n+1}^{m-1}F_k(x)\bigl(g_k(x)-g_{k+1}(x)\bigr).
	\]
	
	По выбору $N$ имеем $\abs{F_k(x)}<\eps$. Поэтому
	\[
	\abs{\sum_{k=n+1}^{m}f_k(x)g_k(x)}
	\le
	\eps\abs{g_m(x)}
	+
	\eps\sum_{k=n+1}^{m-1}\abs{g_k(x)-g_{k+1}(x)}.
	\]
	Так как $\{g_k(x)\}$ монотонна и $\abs{g_k(x)}\le M$, то
	\[
	\sum_{k=n+1}^{m-1}\abs{g_k(x)-g_{k+1}(x)}
	=
	\abs{g_{n+1}(x)-g_m(x)}
	\le 2M.
	\]
	Следовательно,
	\[
	\abs{\sum_{k=n+1}^{m}f_k(x)g_k(x)}
	\le
	M\eps+2M\eps
	=
	3M\eps.
	\]
	Оценка не зависит от $x$, значит хвосты ряда равномерно малы. По критерию Коши ряд $\sum f_kg_k$ сходится равномерно.
\end{proof}

\begin{Theorem}{Признак Дирихле}{}
	Пусть $f_k,g_k\colon\Omega\to\C$. Предположим, что:
	
	\begin{enumerate}
		\item частичные суммы ряда $\sum f_k$ равномерно ограничены: существует $M>0$ такое, что
		\[
		\abs{\sum_{k=1}^{n}f_k(x)}\le M
		\quad \text{для всех } n \text{ и } x\in\Omega;
		\]
		\item $g_k\to0$ равномерно на $\Omega$;
		\item для каждого $x\in\Omega$ последовательность $\{g_k(x)\}$ монотонна.
	\end{enumerate}
	
	Тогда ряд $\sum_{k=1}^{\infty}f_k(x)g_k(x)$ сходится равномерно на $\Omega$.
\end{Theorem}

\begin{proof}
	Обозначим
	\[
	F_j(x)=\sum_{k=n+1}^{j}f_k(x).
	\]
	Из равномерной ограниченности частичных сумм $\sum_{k=1}^{j}f_k(x)$ следует, что $\abs{F_j(x)}\le 2M$ для всех $j$ и $x$.
	
	По преобразованию Абеля:
	\[
	\sum_{k=n+1}^{m}f_k(x)g_k(x)
	=
	F_m(x)g_m(x)
	+
	\sum_{k=n+1}^{m-1}F_k(x)\bigl(g_k(x)-g_{k+1}(x)\bigr).
	\]
	Значит,
	\[
	\abs{\sum_{k=n+1}^{m}f_k(x)g_k(x)}
	\le
	2M\abs{g_m(x)}
	+
	2M\sum_{k=n+1}^{m-1}\abs{g_k(x)-g_{k+1}(x)}.
	\]
	
	Так как $g_k(x)$ монотонна по $k$, сумма модулей разностей телескопируется:
	\[
	\sum_{k=n+1}^{m-1}\abs{g_k(x)-g_{k+1}(x)}
	=
	\abs{g_{n+1}(x)-g_m(x)}
	\le
	\abs{g_{n+1}(x)}+\abs{g_m(x)}.
	\]
	Тогда
	\[
	\abs{\sum_{k=n+1}^{m}f_k(x)g_k(x)}
	\le
	2M\abs{g_{n+1}(x)}+4M\abs{g_m(x)}.
	\]
	Поскольку $g_k\to0$ равномерно, правая часть стремится к нулю равномерно по $x$ при $n,m\to\infty$. По критерию Коши ряд $\sum f_kg_k$ сходится равномерно.
\end{proof}

\subsection{Теорема Дини}

\begin{Theorem}{Теорема Дини}{}
	Пусть $I=[a,b]$, функции $f_n\colon I\to\R$ непрерывны, $f_n(x)\ge0$ для всех $x\in I$, и для каждого $x\in I$ последовательность $\{f_n(x)\}$ монотонно убывает к нулю:
	\[
	f_n(x)\ge f_{n+1}(x),
	\qquad
	f_n(x)\to0.
	\]
	Тогда $f_n\to0$ равномерно на $I$.
\end{Theorem}

\begin{proof}
	Пусть $\eps>0$. Для каждой точки $x\in I$ из поточечной сходимости $f_n(x)\to0$ существует номер $N(x)$ такой, что
	\[
	f_{N(x)}(x)<\eps.
	\]
	Так как $f_{N(x)}$ непрерывна в точке $x$, существует окрестность $U_x$ точки $x$ такая, что для всех $y\in U_x\cap I$:
	\[
	f_{N(x)}(y)<2\eps.
	\]
	
	Семейство $\{U_x\}_{x\in I}$ покрывает компакт $I$. Поэтому существует конечное подпокрытие:
	\[
	I\subset U_{x_1}\cup\cdots\cup U_{x_M}.
	\]
	Положим
	\[
	\widetilde N=\max\{N(x_1),\ldots,N(x_M)\}.
	\]
	
	Возьмём произвольные $y\in I$ и $n>\widetilde N$. Тогда $y\in U_{x_j}$ для некоторого $j$. По монотонности:
	\[
	0\le f_n(y)\le f_{N(x_j)}(y)<2\eps.
	\]
	Значит, $\sup_{y\in I}f_n(y)<2\eps$ при всех $n>\widetilde N$. Следовательно, $f_n\to0$ равномерно на $I$.
\end{proof}

\begin{Corollary}{Следствие для рядов с неотрицательными членами}{}
	Пусть $I=[a,b]$, функции $f_k\colon I\to\R$ непрерывны и $f_k(x)\ge0$ для всех $x\in I$. Если ряд $\sum_{k=1}^{\infty}f_k(x)$ сходится поточечно на $I$ к непрерывной функции $f$, то ряд $\sum_{k=1}^{\infty}f_k$ сходится равномерно на $I$.
\end{Corollary}

\begin{proof}
	Обозначим частичные суммы:
	\[
	S_n(x)=\sum_{k=1}^{n}f_k(x).
	\]
	Так как $f_k(x)\ge0$, то $S_n(x)\le S_{n+1}(x)$, а по условию $S_n(x)\to f(x)$ поточечно.
	
	Рассмотрим остатки $r_n(x)=f(x)-S_n(x)$. Тогда $r_n$ непрерывны, $r_n(x)\ge0$, $r_n(x)\ge r_{n+1}(x)$ и $r_n(x)\to0$ поточечно. По теореме Дини $r_n\to0$ равномерно, то есть
	\[
	\sup_{x\in I}\abs{f(x)-S_n(x)}\to0.
	\]
	Значит, $S_n\to f$ равномерно, поэтому ряд $\sum f_k$ сходится равномерно.
\end{proof}

\subsection{Почленное интегрирование функциональных рядов}

\begin{Theorem}{Почленное интегрирование равномерно сходящегося ряда}{}
	Пусть $I=[a,b]$, функции $f_k\colon I\to\R$ интегрируемы по Риману на $I$, и ряд
	\[
	\sum_{k=1}^{\infty} f_k(x)
	\]
	равномерно сходится на $I$ к функции $f$. Тогда $f$ интегрируема по Риману на $I$ и
	\[
	\int_a^b f(x)\,dx
	=
	\sum_{k=1}^{\infty}\int_a^b f_k(x)\,dx.
	\]
\end{Theorem}

\begin{proof}
	Обозначим частичные суммы и остатки:
	\[
	S_n(x)=\sum_{k=1}^{n}f_k(x),
	\qquad
	r_n(x)=f(x)-S_n(x)=\sum_{k=n+1}^{\infty}f_k(x).
	\]
	Из равномерной сходимости ряда следует, что $r_n\to0$ равномерно на $I$.
	
	Сначала докажем, что $f$ интегрируема по Риману. Для функции $h$ и отрезка $J\subset I$ обозначим её колебание на $J$ через
	\[
	\omega(h;J)=\sup_{x\in J}h(x)-\inf_{x\in J}h(x).
	\]
	По критерию Римана функция интегрируема тогда и только тогда, когда для любого $\eps>0$ существует разбиение $\tau=\{x_0=a,x_1,\ldots,x_N=b\}$ такое, что
	\[
	\sum_{m=0}^{N-1}\omega(h;[x_m,x_{m+1}])(x_{m+1}-x_m)<\eps.
	\]
	
	Зафиксируем $\eps>0$. Выберем $n_0$ так, что
	\[
	\norma{r_{n_0}}_{\infty}<\frac{\eps}{4(b-a)}.
	\]
	Так как $S_{n_0}$ --- конечная сумма интегрируемых функций, она интегрируема. Поэтому существует разбиение $\tau=\{x_0,\ldots,x_N\}$ такое, что
	\[
	\sum_{m=0}^{N-1}
	\omega(S_{n_0};[x_m,x_{m+1}])(x_{m+1}-x_m)
	<\frac{\eps}{2}.
	\]
	На каждом отрезке $[x_m,x_{m+1}]$ имеем $f=S_{n_0}+r_{n_0}$, поэтому
	\[
	\omega(f;[x_m,x_{m+1}])
	\le
	\omega(S_{n_0};[x_m,x_{m+1}])
	+
	\omega(r_{n_0};[x_m,x_{m+1}]).
	\]
	Кроме того,
	\[
	\omega(r_{n_0};[x_m,x_{m+1}])
	\le 2\norma{r_{n_0}}_{\infty}.
	\]
	Следовательно,
	\[
	\sum_{m=0}^{N-1}\omega(f;[x_m,x_{m+1}])(x_{m+1}-x_m)
	<
	\frac{\eps}{2}
	+
	2\norma{r_{n_0}}_{\infty}(b-a)
	<\eps.
	\]
	Значит, $f$ интегрируема по Риману на $I$.
	
	Теперь докажем формулу для интеграла. Так как $f=S_n+r_n$, то
	\[
	\int_a^b f(x)\,dx
	=
	\int_a^b S_n(x)\,dx
	+
	\int_a^b r_n(x)\,dx.
	\]
	При этом
	\[
	\left|\int_a^b r_n(x)\,dx\right|
	\le
	(b-a)\norma{r_n}_{\infty}
	\to0.
	\]
	Значит,
	\[
	\int_a^b f(x)\,dx
	=
	\lim_{n\to\infty}\int_a^b S_n(x)\,dx.
	\]
	Но
	\[
	\int_a^b S_n(x)\,dx
	=
	\int_a^b \sum_{k=1}^{n}f_k(x)\,dx
	=
	\sum_{k=1}^{n}\int_a^b f_k(x)\,dx.
	\]
	Переходя к пределу, получаем
	\[
	\int_a^b f(x)\,dx
	=
	\sum_{k=1}^{\infty}\int_a^b f_k(x)\,dx.
	\]
\end{proof}

\subsection{Почленное дифференцирование функциональных рядов}

\begin{Theorem}{Почленное дифференцирование}{}
	Пусть $I=[a,b]$, функции $f_k\colon I\to\R$ непрерывно дифференцируемы на $I$. Предположим, что:
	
	\begin{enumerate}
		\item существует точка $x_0\in I$, в которой числовой ряд $\sum_{k=1}^{\infty}f_k(x_0)$ сходится;
		\item ряд производных $\sum_{k=1}^{\infty}f_k'$ равномерно сходится на $I$ к функции $g$.
	\end{enumerate}
	
	Тогда ряд $\sum_{k=1}^{\infty}f_k$ равномерно сходится на $I$ к некоторой функции $f$, причём $f$ дифференцируема на $I$ и
	\[
	f'(x)=g(x)=\sum_{k=1}^{\infty}f_k'(x).
	\]
	Иными словами,
	\[
	\left(\sum_{k=1}^{\infty}f_k(x)\right)'
	=
	\sum_{k=1}^{\infty}f_k'(x).
	\]
\end{Theorem}

\begin{proof}
	Обозначим
	\[
	S_n(x)=\sum_{k=1}^{n}f_k(x),
	\qquad
	G_n(x)=\sum_{k=1}^{n}f_k'(x).
	\]
	По условию $G_n\to g$ равномерно на $I$.
	
	Сначала докажем равномерную сходимость ряда $\sum f_k$. Для $m>n$ и $x\in I$ имеем
	\[
	S_m(x)-S_n(x)
	=
	S_m(x_0)-S_n(x_0)
	+
	\int_{x_0}^{x}\bigl(G_m(t)-G_n(t)\bigr)\,dt.
	\]
	Отсюда
	\[
	\abs{S_m(x)-S_n(x)}
	\le
	\abs{S_m(x_0)-S_n(x_0)}
	+
	(b-a)\norma{G_m-G_n}_{\infty}.
	\]
	Ряд $\sum f_k(x_0)$ сходится, поэтому его частичные суммы фундаментальны. Последовательность $G_n$ фундаментальна в равномерной норме, поскольку ряд $\sum f_k'$ равномерно сходится. Значит,
	\[
	\sup_{x\in I}\abs{S_m(x)-S_n(x)}\to0
	\quad \text{при } m,n\to\infty.
	\]
	По критерию Коши $S_n$ равномерно сходится на $I$. Обозначим предел через $f$.
	
	Теперь найдём вид функции $f$. Для каждого $n$ по формуле Ньютона--Лейбница:
	\[
	S_n(x)
	=
	S_n(x_0)+\int_{x_0}^{x}G_n(t)\,dt.
	\]
	Переходя к пределу при $n\to\infty$, получаем
	\[
	f(x)
	=
	\sum_{k=1}^{\infty}f_k(x_0)
	+
	\int_{x_0}^{x}g(t)\,dt.
	\]
	Здесь переход под знак интеграла корректен, потому что $G_n\to g$ равномерно:
	\[
	\left|
	\int_{x_0}^{x}(G_n(t)-g(t))\,dt
	\right|
	\le
	(b-a)\norma{G_n-g}_{\infty}
	\to0.
	\]
	
	Так как $g$ является равномерным пределом непрерывных функций $G_n$, функция $g$ непрерывна. Поэтому по формуле Ньютона--Лейбница функция
	\[
	x\mapsto \int_{x_0}^{x}g(t)\,dt
	\]
	дифференцируема, и её производная равна $g(x)$. Следовательно,
	\[
	f'(x)=g(x)=\sum_{k=1}^{\infty}f_k'(x).
	\]
\end{proof}

\begin{Remark}{Почему нужна сходимость в одной точке}{}
	Равномерной сходимости ряда производных недостаточно, чтобы восстановить сами функции $f_k$ с точностью до констант. Условие сходимости $\sum f_k(x_0)$ в одной точке фиксирует эти константы.
\end{Remark}